\subsection{Filtering}
\label{SEC:Filtering}
\key*{}[filtering]

%%%%%%%%%%%%%%%%%%%%%%%%%%%%%%%%%%%%%%%%%%%%%%%%%%%%%%%%%%%%%%%%%%%%%%%%%%%%%%%%

In spectral-element methods, the solution within an individual element
can be transformed into a hierarchical set of modal coefficients which
are representative of the energy contained in each mode.
%
A well-resolved numerical solution will have minimal energy contained
within the higher modes, and the modal coefficients should have an
exponential decay rate from lowest to highest modes.
%\cite{21, 22, 32, 33}%% From AIAA-2022-1813
%
For this reason, an exponential filter is used as a source of
numerical dissipation for stabilizing the simulations.

%%%%%%%%%%%%%%%%%%%%%%%%%%%%%%%%%%%%%%%%%%%%%%%%%%%%%%%%%%%%%%%%%%%%%%%%%%%%%%%%
%% Note for future reference:
%%
%% For the diagonal filter matrix sigma(m) in GFR,
%% see: Spiegel AIAA-2022-1813 Eq.5,
%% but that index goes from m=0:N_M "for the N_M polynomial modes (p.2)".
%% The polynomial order is degree P (p.2) and 
%% the number of solution points is N_P=P+1.
%% His previous paper numbers the modal coefs 0:P and the nodal coefs 1:P=1.
%% Since a polynomial of degree P has terms x^0 + x^1 + ... x^P
%% it will have just as many modes 0:P with 0 being the constant mode.
%% In this context, it seems N_M=P even though it is not explicitly defined.
%% However, in practice the arrays in GFR begin at index 1.
%% 
%% Hesthaven & Warburton (2010) \cite{Hesthaven_2010}
%% p.20:
%%   N=polynomial order; N_P=polynomial points; so N_p=N+1.
%%   The modal coefs are 1:N_P and the nodal coefs are 1:N_P too.
%% p.130:
%%   Defines sigma(eta) with eta=(n-1)/N; n=1...N_p;
%%   Their Matlab code on that page also provides a
%%     1-based diagonal filter matrix.
%%
%% Sheshadri (2016) PhD \cite{sheshadri16}
%% p. 11:  The nodal coefs are n=0:P.
%% p. 99:  The modal coefs are n=0:P too.
%% p.107:
%%   Defines sigma(eta) with eta=n/P; n=0...P;
%%
%% We use Hesthaven's notation here, because it is consistent with the
%% 1:P+1 modal form of the polynomial used in our introduction.

This diagonal exponential filter matrix is given by\cite{Hesthaven_2010}
\begin{align}
    \label{EQN:ExpFilterMatrix}
    \sigma (n)
 =& \begin{cases}
       1
    &  0 \leq \eta \leq \eta_c
    \\
       \exp \left[ -\alpha \left( \frac{\eta-\eta_c}{1-\eta_c} \right)^s
            \right]
    &  \eta_c \leq \eta \leq 1
    \end{cases}
  &
    \begin{aligned}
       \eta   = & \;\text{modal fraction}, (n-1) / \PO
    \\ \eta_c = & \;\text{modal fraction cutoff}
    \\ \alpha = & \;\text{filter strength}
    \\ s      = & \;\text{filter order}
    \end{aligned}
\end{align}
%
where $n=1,\ldots,\PO+1$.

%%%%%%%%%%%%%%%%%%%%%%%%%%%%%%%%%%%%%%%%%%%%%%%%%%%%%%%%%%%%%%%%%%%%%%%%%%%%%%%%

%Three user-specified parameters control the aggressiveness of the filter,
%and their effects are illustrated in \fref{FIG:Filter}.
%%
%%
%(1) The parameter $\eta_c$ specifies the modal fraction cutoff above which
%filtering will be applied.
%%
%Everything below that modal fraction cutoff will be entirely preserved.
%%
%Most importantly, the lowest mode
%(which defines the average value in the cell)
%remains unchanged in order to ensure conservation of
%mass, momentum, and energy.
%%
%%
%(2) The parameter $\alpha$ specifies the strength of the filter.
%As $\alpha$ increases, the filter value at the highest mode(s)
%becomes smaller, which means that more damping is applied.
%%
%%
%(3) Lastly, the parameter $s$ specifies the filter order
%which controls the exponential decay rate
%(\thatis how parabolic the filter function is on a log plot).
%As $s$ increases, filtering has substantially less of an effect on
%the lower modes.

%%%%%%%%%%%%%%%%%%%%%%%%%%%%%%%%%%%%%%%%%%%%%%%%%%%%%%%%%%%%%%%%%%%%%%%%%%%%%%%%

Three user-specified parameters, $[\eta_c, \alpha, s]$, control the
aggressiveness of the filter, and their effects are illustrated in
\fref{FIG:Filter}.\\[12pt]

%%%%%%%%%%%%%%%%%%%%%%%%%%%%%%%%%%%%%%%%%%%%%%%%%%%%%%%%%%%%%%%%%%%%%%%%%%%%%%%%

\begin{figure}[htb]
  \begin{centering}
    \begin{subfigure}{2.0in}
      \insertimage[width=\textwidth]{Filter/Filter_etac}
      \caption{Effect of $\eta_c$}
      \label{FIG:Filter_etac}
    \end{subfigure}%
    \begin{subfigure}{2.0in}
      \insertimage[width=\textwidth]{Filter/Filter_alpha}
      \caption{Effect of $\alpha$}
      \label{FIG:Filter_alpha}
    \end{subfigure}%
    \begin{subfigure}{2.0in}
      \insertimage[width=\textwidth]{Filter/Filter_s}
      \caption{Effect of $s$}
      \label{FIG:Filter_s}
    \end{subfigure}%
    %%%%%%%%%%%%%%%%%%%%%%%%%%%%%%%%%%%%%%%%%%%%%%%%%%%%%
    \caption{Exponential filter function, $\sigma$.}
    \label{FIG:Filter}
  \end{centering}
\end{figure}

%%%%%%%%%%%%%%%%%%%%%%%%%%%%%%%%%%%%%%%%%%%%%%%%%%%%%%%%%%%%%%%%%%%%%%%%%%%%%%%%
\begin{tabular}{%
      @{\hspace{2.0em}} l%
      @{\hspace{1.0ex}} c%
      @{\quad} >{\raggedright\arraybackslash}p{0.8\linewidth}%
      @{}}
  %
  $\eta_c$ & $\coloneq$ &
  This parameter specifies the modal fraction cutoff above which filtering will
  be applied.
  %
  Everything below that modal fraction cutoff will be entirely preserved.
  %
  Most importantly, the lowest mode (which defines the average value in the
  cell) remains unchanged in order to ensure conservation of mass, momentum, and
  energy.
  %
  Valid input values for this parameter are real values on the interval
  $[0.0,1.0]$, although a value of $1.0$ is the same as using no filter at all
  because it will preserve all modes of the solution.
  %
  A value of $0.0$ will apply some amount of filtering to all modes except the
  lowest\slash constant mode.\\
  %
  & & \\
  %   
  $\alpha$ & $\coloneq$ & 
  This parameter specifies the strength of the filter, \thatis the highest mode
  will be multiplied by the value $\exp\left(-\alpha\right)$.
  %
  As $\alpha$ increases, the filter value at the highest mode(s) becomes
  smaller, which means that more damping is applied.
  %
  Valid input values for this parameter are real values between $0.0$ and
  roughly $36.0$, where a value of $36.0$ is nearly equivalent to multiplying
  the highest mode by 64-bit machine epsilon, \thatis
  $\ln\left(2^{-52}\right)\approx-36.0$.\\
  %
  & & \\
  %
  $s$ & $\coloneq$ & 
  This parameter specifies the filter order which controls the exponential decay
  rate (\thatis how parabolic the filter function is on a log plot).
  %
  As $s$ increases, filtering has substantially less of an effect on the lower
  modes.
  %
  Valid input values for this parameter are positive integers greater than or
  equal to 2; however, anything greater than 32 will essentially result in a
  filter that multiplies the highest mode by $\exp\left(-\alpha\right)$ applies
  negligible filtering to any lower modes that are larger than the cutoff mode.
  %
  In practice, commonly used values are $[2,4,8,12,16,32]$.
  %
\end{tabular}
\\[12pt]

%%%%%%%%%%%%%%%%%%%%%%%%%%%%%%%%%%%%%%%%%%%%%%%%%%%%%%%%%%%%%%%%%%%%%%%%%%%%%%%%

GFR uses a three-pass filtering approach.
%
First, a \key*{\emph{sponge filter}}[filtering!sponge filter] can be applied
to the downstream region starting from a user-specified x-location.
This filter region is often used to damp pressure fluctuations so they do not
reflect from the outflow boundary and contaminate the upstream solution.
%
Second, a {\key*{\emph{local filter}}[filtering!local filter] can be applied
to those cells identified by the
\key[key:marker]{troubled cell marker}[marker] in \sref{SEC:Markers}.
%
Usually these cells require some significant level of filtering
due to a shock or some local instability.
%
Third, a weak \key*{\emph{global filter}}[filtering!global filter] can be
applied to all cells in the domain to remove any high-frequency noise in the
solution; NOTE, use of the global filter should be avoided if possible since it
filters \emph{everything} including cells that do not need filtering.

%%%%%%%%%%%%%%%%%%%%%%%%%%%%%%%%%%%%%%%%%%%%%%%%%%%%%%%%%%%%%%%%%%%%%%%%%%%%%%%%

It is recommended to check for troubled cells and perform 
filtering operations after each Runge-Kutta subiteration.
%
This is the highest filtering frequency possible, with other researchers
opting to filter only after the final Runge-Kutta stage or after a
certain number of time steps.
%
Filtering at such a high frequency allows for the use of a larger
time step, but it can reduce the accuracy of the solution.
%
This decrease in accuracy is usually outweighed by the increased
time step size, which is often the most limiting computational factor
in time-to-solution performance.

%%%%%%%%%%%%%%%%%%%%%%%%%%%%%%%%%%%%%%%%%%%%%%%%%%%%%%%%%%%%%%%%%%%%%%%%%%%%%%%%

\SeeAlso{}
\noindent%
\begin{tabular}{@{\qquad} lll @{}}
\begin{varwidth}[t]{0.3\textwidth}
\inp[nmlLimiter::GlobalFilterInterval]{Global\_Filter\_Interval}, \\
\inp[nmlLimiter::GlobalFilterOption]{Global\_Filter\_Option}, \\
\inp[nmlLimiter::GlobalFilterOrder]{Global\_Filter\_Order}, \\
\inp[nmlLimiter::GlobalFilterAlpha]{Global\_Filter\_Alpha}, \\
\inp[nmlLimiter::GlobalFilterEtac]{Global\_Filter\_Etac}, \\
%
\end{varwidth}
&
\begin{varwidth}[t]{0.3\textwidth}
\inp[nmlLimiter::LocalLimiterInterval]{Local\_Limiter\_Interval}, \\
\inp[nmlLimiter::LocalFilterOrder]{Local\_Filter\_Order}, \\
\inp[nmlLimiter::LocalFilterAlpha]{Local\_Filter\_Alpha}, \\
\inp[nmlLimiter::LocalFilterEtac]{Local\_Filter\_Etac}, \\
%
\end{varwidth}
&
\begin{varwidth}[t]{0.3\textwidth}
\inp[nmlLimiter::SpongeFilterInterval]{Sponge\_Filter\_Interval}, \\
\inp[nmlLimiter::SpongeFilterOption]{Sponge\_Filter\_Option}, \\
\inp[nmlLimiter::SpongeFilterOrder]{Sponge\_Filter\_Order}, \\
\inp[nmlLimiter::SpongeFilterAlpha]{Sponge\_Filter\_Alpha}, \\
\inp[nmlLimiter::SpongeFilterEtac]{Sponge\_Filter\_Etac}, \\
\inp[nmlLimiter::SpongeFilterXStart]{Sponge\_Filter\_X\_Start}\end{varwidth}
\end{tabular}

%%%%%%%%%%%%%%%%%%%%%%%%%%%%%%%%%%%%%%%%%%%%%%%%%%%%%%%%%%%%%%%%%%%%%%%%%%%%%%%%

